\section{The first variation on the local curve}\label{sec:original}
\begin{proposition}\label{prop:original}
The original index satisfies \eqref{eq:first} on every transpose orbit.
\end{proposition}
\subsection{The fixed virtual zero-cycle}


Fix a theta characteristic $S$ as reference target, $S^2=K$. For the original antidiagonal action, tensor the weight-$j$ piece of a fixed pair at $S$ by $(L/S)^j$ to obtain the full Picard family. In each weight-$j$ summand the inverse twist is $(S/L)^j$. Both operations preserve the section and coordinate maps, commute with base change, and preserve Euler characteristic since $L/S$ has degree zero on the curve. They identify the family of fixed schemes with the product over $B$. In a local frame of $L/S$, each weight-zero map of the trace-free pair complex has total twist zero. The relative Atiyah class in the pair-parameter directions holds $L/S$ fixed. Thus the fixed obstruction morphism and virtual sheaf are pulled back from the reference scheme \cite[\S2]{W}. We use the residual volume action before specialization only to perform equivariant pushforward and form the moving Euler ratios.

At the function field of $C$ a fixed pair is a finite cyclic homogeneous quotient of the two-variable polynomial algebra supported at the origin. Every weight space of that algebra is a monomial times the one-variable ring in the product of coordinates. Its homogeneous ideal in that space is generated by a power of this product, since the product acts nilpotently. Hence the generic quotient is monomial. Its diagram $\lambda$ is determined by the ranks of the weight bundles $E_j$ on $C$: on each diagonal of the diagram the occupied boxes form an initial segment. These ranks give a decomposition into open and closed loci $F_{\lambda,n}$ at Euler grading $n$. Swapping the two copies of $S$ carries $F_{\lambda,n}$ to $F_{\lambda^T,n}$.

On such a locus the fixed tangent at volume one is $P_0-P_0^\vee$ for the zero-weight part of the polarization $P$ of Lemma~\ref{lem:root}. Its virtual dimension is zero. Its ordinary virtual structure sheaf belongs integrally to the dimension-zero filtration of $G_0(F_{\lambda,n})$. We recall the integral filtration argument. For a vector bundle $p:V\to X$ of rank $r$, flat pullback carries $F_jG_0(X)$ onto $F_{j+r}G_0(V)$. Indeed Chow groups map onto the associated graded groups for the dimension filtration by taking structure sheaves, compatibly with flat pullback, and integral Chow homotopy invariance identifies $CH_j(X)$ with $CH_{j+r}(V)$. Subtract a lift in highest degree and descend in the filtration. In particular the pieces below dimension $r$ on $V$ vanish. Homotopy invariance of $G_0$ makes the inverse the zero-section Gysin map. For a global two-term bundle resolution of the fixed perfect obstruction theory (available on the quasi-projective scheme), the lifted intrinsic normal cone in the obstruction bundle of rank $r$ has dimension $r+\mathrm{vdim}$. Its structure sheaf pulled back along the zero section gives the assertion. This argument uses ordinary coherent $K$-theory after forgetting the residual torus \cite[\S3]{FG}.

Write $e_{\lambda,n}$ for this class with the counting sign, the moving spin sign and the fixed determinant line at volume one included. It is an integral sum of closed point classes. The line does not affect a point class. The sum over disconnected components retains their respective signs. Transposition preserves $e$ under transport. In fact the two polarizations differ by $V+V^\vee$ at volume one by curve Serre duality, with $V$ the push of the coordinate-weighted endomorphism bundle up to sign. This leaves the determinant of the fixed part and parity of the moving rank unchanged. The swap transports the fixed stable-pair obstruction theory itself. The degree of $e_{\lambda,n}$ is $[q^n]Z_\lambda(q)$. Here the numerical input is available \emph{partition by partition}: combining the three antidiagonal universal series of Monavari at balanced levels cancels the factor $F_\lambda$ and gives the hook product with the signed counting convention of \cite{W}. The generic diagram of each full-torus fixed component in that computation is exactly the one determined by the ranks above. This uses a numerical calculation only for the degree of an already integral zero-cycle class \cite[Theorem~1.7 and \S9]{M}.

\subsection{Positive exchanged grading}

A nonzero moving weight $t^j V$ in $P$ has spin-normalized Euler ratio $\lambda_{-\kappa^{-1} t^j}(V)/\lambda_{-t^j}(V)$, where $\lambda_z(V)=\sum_{d\ge0}z^d[\bigwedge^dV]$, after extraction of signs and scalar characters. At negative $j$ reflect the ratio to a positive power with dual coefficient and inverted volume parameter. These identities hold equivariantly on the proper fixed scheme. Every nonconstant coefficient of a normalized ratio is divisible by $1-\kappa$ by the splitting principle, and its constant term at $t=0$ is a volume unit. Divide the coefficient at $t^k$, $k>0$, by $1-\kappa$ and then specialize. Only one moving ratio contributes. A positive $j\mid k$ supplies $-\psi^{k/j} V$; a negative $-j$ with $j\mid k$ supplies $+\psi^{k/j} V^\vee$. All other factors specialize to their zeroth order, giving $e_{\lambda,n}$. This order of operations requires neither a volume derivative of a residual equivariant sheaf nor division by two. The actual graded family insertions can be restricted to closed points of the reference fixed scheme before pushforward by the dimension filtration argument above.

On the minimal quotient the plane Hilbert tangent gives the positive-half monomials $N^{a_x+1}P_2^{-\ell_x}$, where $N=tL$, $P_2=K/(tL)$ and $a_x=h_x-1-\ell_x$. The polarized tangent at the minimum is their curve push modulo hyperbolic summands, by the following comparison. The Hilbert tangent at a monomial ideal has pairs $w,NP_2/w$. Their curve pushforwards are Serre dual with opposite signs. The difference between the positive-half polarization and the Koszul polarization before pushforward is antisymmetric under $w\mapsto NP_2/w$. There is no fixed monomial in the two independent normal variables, so the difference is a sum of opposite pairs. They push to hyperbolic summands. A hyperbolic summand $A+A^\vee$ contributes zero to the positive-$t$ derivative. The resulting minimal insertion is the negative inner sum of Theorem~\ref{thm:first}.

To compute its change at an arbitrary closed fixed pair, the section and its monomials embed the minimal quotient bundle $E^{\min}$ into $E$. Monomials outside the diagram vanish generically and hence everywhere by torsion-freeness of $E$. The map on the remaining bundle is generically an isomorphism, so its kernel is zero. This also applies when a weight has rank greater than one. For a curve point $p$ let $d_j(p)$ be the length of the weight-$j$ part of the finite cokernel there. Write
\[
V(z)=\sum_j\operatorname{rk}(E_j)z^j,\qquad Q(z)=\sum_j d_j(p)z^j,\qquad f^*(z)=f(z^{-1}).
\]
Using $z=tL_p$ in the family $K$-group, the change in the Koszul polarization from this support point is
\[
F=Q-(1-z)(QV^*-VQ^*).
\]
Indeed in $K^0(C)$ finite-length duals acquire a minus sign, a point squared is zero, and multiplication by an ordinary bundle supported against a point sees its rank. The constant lines $S_p,K_p$ have class one. All family weights are precisely the powers of $L_p$. Thus the insertion change is the sum, for $j>0$ dividing $k$, of $[z^j]\bigl(F^*-F\bigr)$ times $L_p^k$. Explicitly
\[
F^*-F=Q^*-Q+(2-z-z^{-1})(QV^*-VQ^*).
\]
It changes sign on simultaneously inverting $Q,V$. This proves cancellation between distinct transposed shapes with their transported zero-cycles.

For a self-transposed shape $V=V^*$, the right hand side is $(Q^*-Q)(1-(2-z-z^{-1})V)$. Each insertion is now an integer combination of differences $d_j(p)-d_{-j}(p)$. One must justify cancellation without choosing a transposition-invariant representative of a zero-cycle. On the reference fixed scheme form the perfect insertions
\[
J_{j,k}=R\rho_*[(E_j-E^{\min}_j)L^k],
\]
using the reference weight bundles on $C$, with their torus weights omitted, and pulling $L$ from $C\times B$. At each fixed pair their restrictions are $\sum_p d_j(p)L_p^k$. The construction needs the minimal map only at closed points; the class in brackets is globally defined by a difference of bundles. Swapping the coordinates on the reference target carries $J_{j,k}$ to $J_{-j,k}$, keeping the Picard factor fixed. Their difference has zero push against the invariant class $e_{\lambda,n}$ by the projection formula and the change of variables under the swap. For distinct transposes the same global insertions can be used on transported representatives. This argument cancels actual differences in integral $K^0(B)$, including any possible additive torsion.

What remains on each orbit is therefore the minimal insertion on each shape multiplied by the degree of its fixed zero-cycle. Summation in $n$ proves the formula.
